\begin{enumerate}
    \item Donner la définition du volume molaire et de la mole.
    \item On considère un échantillon de fer \ce{Fe} de masse
          $m=\qty{5.6}{\g}$.
          \begin{enumerate}
              \item Calculer la quantité de matière contenue dans cette masse de
                    fer.
              \item Déterminer le nombre d'atomes contenus dans cet échantillon.
          \end{enumerate}
    \item Un flacon contient un volume $V=\qty{230}{\cubic\cm}$ d'éthanol
          \ce{C2H6O} pur à l'état liquide dont la densité par rapport à l'eau
          $d=0,79$.
          \begin{enumerate}
              \item Calculer la quantité de matière d'éthanol contenue dans ce
                    flacon.
              \item En déduire la masse de cette quantité d'éthanol.
          \end{enumerate}
    \item Une bouteille contient un volume $V=\qty{230}{\cubic\cm}$ du dioxygène
          \ce{O2} gazeux sous la pression $P=\qty{1033}{\hPa}$
          et à la température $\theta=\qty{25}{\degreeCelsius}$.
          \begin{enumerate}
              %\item Déterminer la densité du dioxygène par rapport à l'air.
              \item Calculer la quantité de matière du gaz dioxygène qui se
                    trouve dans cette bouteille (en le considérant comme un gaz
                    parfait).
              \item Déterminer la valeur du volume molaire dans les conditions
                    précédentes.
              \item Quelle est la pression qu'on doit exercer sur l'échantillon
                    du gaz précédent à la température
                    $\theta^{\prime}=\qty{20}{\degreeCelsius}$ pour que son
                    volume devient $V^{\prime}=\qty{0,8}{\liter}$.
          \end{enumerate}
\end{enumerate}
\begin{donnees}
    \begin{itemize}
        \item $M(\ce{C2H6O})=\qty{46}{\gram\per\mol}$, \;
              $M(\ce{Fe})=\qty{56}{\gram\per\mol}$, \;
              $M(\ce{O2})=\qty{32}{\gram\per\mol}$.
        \item $\rho_{\text{eau}}=\qty{1}{\gram\per\cubic\cm}$, \;
              $N_{A}=\qty{6.02e23}{\per\mol}$, \;
              $R=\qty{8,314}{\J\per\mol\per\K}$,
        \item $\qty{1}{\liter}=\qty{1e-3}{\cubic\meter}$, \;
              $\qty{1}{\hecto\Pa}=\qty{100}{\Pa}$.
    \end{itemize}
\end{donnees}
